% Generated draft: requires mathematical and source review.
\documentclass[11pt]{article}
\usepackage[margin=25mm]{geometry}
\usepackage{fontspec}
\setmainfont{Latin Modern Roman}
\usepackage{amsmath,amssymb,mathtools,unicode-math}
\setmathfont{Latin Modern Math}
\usepackage{enumitem,needspace,xurl,hyperref}
\setlength{\parindent}{0pt}
\setlength{\parskip}{5pt}
\setlength{\emergencystretch}{3em}
\newcommand{\N}{\mathbb{N}}
\newcommand{\Z}{\mathbb{Z}}
\newcommand{\Q}{\mathbb{Q}}
\newcommand{\R}{\mathbb{R}}
\newcommand{\C}{\mathbb{C}}
\newcommand{\F}{\mathbb{F}}
\newcommand{\A}{\mathbb{A}}
\newcommand{\PP}{\mathbb{P}}
\newcommand{\E}{\mathbb{E}}
\newcommand{\eps}{\varepsilon}
\newcommand{\dd}{\,\mathrm{d}}
\begin{document}
\section*{UnsolvedMath ID: 30003485}
\textbf{Source problem code:} catalogue entry
\subsection{Definitions and mathematical statement}
The supplied record asks for a classification of phylogenetic network topologies by their distinguishability under three group-based models: the Cavender--Farris--Neyman model (CFN), the Kimura two-parameter model (K2P), and the Kimura three-parameter model (K3P). CFN has a two-state alphabet, conventionally indexed by \(\Z/2\Z\); K2P and K3P have four-state alphabets, conventionally indexed by \((\Z/2\Z)^2\). These model names alone do not specify the stochastic construction on a network.

Let \(X\) be a fixed leaf-label set with \(|X|=n\). Write \(\mathcal C_k(X)\) for the intended class of \(k\)-cycle phylogenetic networks on \(X\), and \(\mathcal L_1(X)\) for the intended class of level-one phylogenetic networks on \(X\). Here interpreting \(n\) as the number of labelled leaves follows the usual phylogenetic reading; the supplied statement does not explicitly define it. The precise network conventions remain unresolved. In a common rooted convention, level one means that each biconnected component of the underlying undirected graph contains at most one reticulation vertex; a single-cycle network has one such cycle, whose length may be denoted by \(k\). The record does not establish that these are its intended conventions, or specify binaryness, directions retained in the topology, suppression of degree-two vertices, or admissible cycle lengths.

To make the possible meaning of distinguishability explicit, suppose a network parameterization has been fixed. For a model \(M\) with state alphabet \(A_M\), let
\[
\phi_{N,M}\colon\Theta_{N,M}\longrightarrow\Delta(A_M^X),\qquad
\Delta(A_M^X)=\left\{(p_a)_{a\in A_M^X}:p_a\geq0,\ \sum_{a\in A_M^X}p_a=1\right\},
\]
where \(\Theta_{N,M}\) is the chosen stochastic parameter domain. Define the attainable distribution set and its complex Zariski closure by
\[
\mathcal P_M(N)=\phi_{N,M}(\Theta_{N,M}),\qquad
V_M(N)=\overline{\mathcal P_M(N)}^{\,\mathrm{Zar}}\subseteq\C^{A_M^X}.
\]
The choice of \(\Theta_{N,M}\), including root distributions, edge constraints, reticulation mixing, and inclusion of boundary parameters, is not supplied. Thus these expressions describe a framework for specifying the problem, rather than a completed choice of network model.

One possible algebraic distinguishability relation is mutual noncontainment:
\[
D_M^{\mathrm{nc}}(N,N')\quad\Longleftrightarrow\quad
V_M(N)\not\subseteq V_M(N')\ \text{and}\ V_M(N')\not\subseteq V_M(N).
\]
Other possible meanings include merely \(V_M(N)\neq V_M(N')\), inequality of attainable stochastic distribution sets, or a specified generic statistical identifiability condition. These criteria must not be silently identified; the supplied record does not select one.

\textbf{Requested classification, subject to these unresolved specifications.} For each fixed admissible \(n\), each admissible cycle length \(k\), and each
\[
M\in\{\mathrm{CFN},\mathrm{K2P},\mathrm{K3P}\},
\]
give necessary and sufficient conditions on a pair \(N,N'\in\mathcal C_k(X)\) for distinguishability under \(M\). Extend the classification to all pairs \(N,N'\in\mathcal L_1(X)\). Equivalently, once an intended relation \(D_M\) and network conventions are fixed, determine exactly the sets
\[
\{(N,N')\in\mathcal C_k(X)^2:D_M(N,N')\}
\quad\text{and}\quad
\{(N,N')\in\mathcal L_1(X)^2:D_M(N,N')\}.
\]
No numerical range for \(n\) or \(k\), or treatment of comparisons between different cycle lengths, is supplied. A unique fully specified formulation therefore requires clarification of the network classes, parameterizations, and distinguishability criterion.
\subsection{Short English statement}
For a fixed number of leaves, classify which pairs of cycle phylogenetic networks can be distinguished under each of CFN, K2P, and K3P, and extend the classification to level-one networks. The source leaves the network conventions, parameter restrictions, and precise meaning of distinguishability unspecified. Its partially solved status and reported special cases are supplied dated curation requiring independent verification.
\subsection{Sources}
\begin{itemize}
\item Supplied local UnsolvedMath registry record, permanent ID \(30003485\), source problem code \texttt{catalogue entry}, title \emph{Distinguishability of Cycle and Level-One Phylogenetic Networks}; collection: \emph{Oberwolfach Reports Open Problems}. No individual proposer is named. The proposed year is recorded as 2017.
\item Supplied workshop-report citation: \emph{Algebraic Statistics (2018)}, Oberwolfach Reports, pp.\ 1207--1279, DOI \texttt{10.4171/owr/2017/20}. \url{https://doi.org/10.4171/owr/2017/20}. The record attributes the problem to the workshop-report collection; it gives no problem-specific page or speaker attribution. The supplied title year 2018 and proposed year 2017 are retained without reconciliation.
\item Supplied literature source: ``Gross and Long on cycle networks.'' \url{https://doi.org/10.1137/17M1134238}.
\item Supplied literature source: ``Hollering and Sullivant on algebraic matroids.'' \url{https://doi.org/10.1016/j.jsc.2020.04.012}.
\item Supplied literature source: ``recent level-one identifiability.'' \url{https://arxiv.org/abs/2607.12919}.
\item Literature assessment dated exactly \texttt{2026-08-21}: the supplied curation labels the problem partially solved and reports known four-leaf CFN cases and generic K2P and K3P large-cycle identifiability, while describing the classification for every cycle pair and all listed models as incomplete. This is supplied dated curation derived from the OWR record, not an independently verified literature assessment. The record explicitly requests independent verification of time-sensitive status and cited claims.
\item The supplied statement status is \texttt{corrected\_verified}; its accompanying explanation says that notation, typography, and sentence structure were repaired while preserving mathematical content. This records an editorial assessment and does not establish mathematical correctness or independent verification here.
\end{itemize}
\end{document}
